\section{Literature Review}
\label{sec:literature}

This section reviews two streams of literature that define the paper's application structure and algorithmic positioning.
The first stream is electric-bus charging and wireless infrastructure planning, which motivates the interaction among wireless facilities, battery choices, and existing plug-in fast-charging resources.
The second stream is stochastic facility-location modeling and decomposition, which provides the methodological context for two-stage mixed-integer planning with scenario-dependent recourse.
The review is organized to keep the transportation discussion concise and to focus on the shared-resource mixed-integer recourse structure used by the proposed algorithm.

\subsection{Electric-bus charging and wireless infrastructure planning}
\label{subsec:lit-wireless-bus}

% TODO: Add verified citations for SWC/DWC planning, electric-bus battery sizing, plug-in charging scheduling, and route-level energy management.

Electric-bus charging studies commonly coordinate infrastructure deployment, battery sizing, vehicle or charging schedules, and service feasibility.
These decisions are often formulated as integrated optimization models because charging resources, route energy requirements, and schedule feasibility interact.
Depot and plug-in fast-charging studies emphasize charger availability and charging schedules, whereas opportunity-charging studies add charging possibilities during service.

Wireless charging infrastructure planning extends this setting by considering stationary wireless charging at stops and dynamic wireless charging on road links.
This literature provides modeling elements that are central to our application: candidate wireless facilities, route-level energy feasibility, battery-capacity implications, and charging-operation constraints.
It also motivates the coexistence of wireless facilities with existing plug-in fast-charging resources.

Our use of this stream is structural.
We adopt wireless-facility placement, route-level battery choices, representative-day charging operations, and existing plug-in fast chargers as the application objects of the model.
We do not focus on real-time charging control, bus dispatch, or traffic operations; instead, these objects define a long-term planning setting whose shared-resource mixed-integer recourse structure is reviewed methodologically next.

\subsection{Stochastic facility-location models and decomposition methods}
\label{subsec:lit-stochastic-decomposition}

% TODO: Add verified citations for stochastic facility location and stochastic network design.
% TODO: Add verified citations for integer L-shaped methods and mixed-integer recourse.
% TODO: Add verified citations for partial Benders or lifted master reformulations.
% TODO: Add verified citations for relaxation-based recourse switching and Benders subproblem decomposition.
% TODO: Add verified citations for logic-based, combinatorial, and monotonicity-based Benders cuts.

Stochastic facility-location and network-design models provide the optimization background for here-and-now infrastructure decisions under uncertain future conditions.
In these models, scenarios typically represent uncertain demand, service requirements, or operating conditions, and the recourse problem evaluates the consequence of a first-stage design after a scenario is realized.
Benders decomposition, integer L-shaped methods, and logic-based Benders decomposition provide the standard decomposition language for such two-stage models, especially when the deterministic equivalent is too large or the recourse problem contains discrete decisions.

Within this decomposition family, computational performance is shaped by two broad design dimensions.
The first is the decomposition or reformulation choice: what information is retained in the master problem, what remains in the scenario subproblems, and whether the subproblems can be decomposed or relaxed without losing useful bound information.
The second is the separation strategy: how feasibility and optimality information from the subproblems is converted into cuts, and whether those cuts exploit problem-specific structure beyond a basic no-good exclusion.

Research on decomposition for mixed-integer recourse has developed techniques along both dimensions.
On the reformulation side, stronger master representations, partial recourse descriptions, relaxed subproblems, and internal subproblem decompositions can improve lower bounds or reduce separation cost.
On the cut side, integer L-shaped cuts, logic-based Benders cuts, combinatorial cuts, valid inequalities, and dominance-based cuts can exploit structure that is not captured by generic linear duality.
The effectiveness and validity of these enhancements are structure dependent, which is why successful decomposition algorithms are often tailored to the recourse model rather than transferred directly from one application to another.

This paper follows that tailored-decomposition perspective.
The electric-bus wireless-charging setting creates a shared-resource mixed-integer recourse structure with scenario-dependent battery choices, route-level operating decisions, and plug-in fast-charger capacity coupling.
The proposed method instantiates the two design dimensions above through a reformulation-enabled integer L-shaped algorithm: it changes the master--subproblem organization, uses a relaxation--verification separation scheme, and derives cuts from the joint facility--battery structure while retaining exact recourse verification.
